\chapter{NGHIÊN CỨU LIÊN QUAN}

\section{Các nghiên cứu liên quan}

Khả năng tính toán của mạng nơ-ron đã được nghiên cứu rộng rãi trong nhiều thập kỷ qua. Các nghiên cứu tiên phong về RNN đã chứng minh rằng RNN có thể mô phỏng máy Turing về mặt lý thuyết \cite{siegelmann_sontag_1995}, nghĩa là chúng có khả năng biểu diễn bất kỳ hàm tính toán được nào. Gần đây hơn, kiến trúc Transformer cũng được chứng minh có khả năng biểu diễn tương đương \cite{perez_2021}, củng cố thêm quan điểm về tính phổ quát của mạng nơ-ron.

Tuy nhiên, điều quan trọng cần phân biệt là tính khả thi về mặt biểu diễn hoàn toàn khác với khả năng học được chính xác thông qua gradient descent. Một mạng nơ-ron có thể có khả năng biểu diễn một hàm phức tạp, nhưng điều đó không đảm bảo rằng thuật toán tối ưu hóa như gradient descent có thể tìm được tập tham số phù hợp để thực hiện hàm đó từ dữ liệu huấn luyện hữu hạn. Đây chính là khoảng cách then chốt mà các nghiên cứu trước đó chưa giải quyết triệt để.

Nhiều kiến trúc chuyên biệt đã được đề xuất để thu hẹp khoảng cách này. Neural GPU mở rộng khả năng của mạng tích chập để xử lý các phép toán số học song song. NALU được thiết kế đặc biệt để học các phép toán số học như cộng, trừ, nhân và chia. Mặc dù các kiến trúc này đạt được kết quả thực nghiệm tốt trong một số trường hợp, chúng thiếu nền tảng lý thuyết vững chắc về việc tại sao và khi nào quá trình học hội tụ đến giải pháp chính xác.

Nghiên cứu đã cho thấy sự khác biệt cơ bản ở chỗ chứng minh lý thuyết chặt chẽ dựa trên khung NTK, đảm bảo rằng mạng nơ-ron không chỉ có khả năng biểu diễn mà thực sự có thể học được các phép toán chính xác thông qua gradient descent với lượng dữ liệu huấn luyện có thể kiểm soát được.

\section{Cơ sở lý thuyết}

\subsection{Ký hiệu và các khái niệm cơ bản}

Trước khi đi vào phân tích chi tiết khung lý thuyết NTK, các ký hiệu toán học cơ bản được sử dụng xuyên suốt báo cáo này là:

\begin{enumerate}
    \item \textbf{Vector và ma trận:}
    \begin{itemize}
        \item Vector cột $n$ chiều: $\mathbf{x} \in \mathbb{R}^n$
        \item Phần tử thứ $i$ của vector: $x_i$ hoặc $[\mathbf{x}]_i$
        \item Ma trận đơn vị kích thước $n \times n$: $I_n$
        \item Tích vô hướng: $\langle \mathbf{x}, \mathbf{y} \rangle = \mathbf{x}^\top \mathbf{y} = \sum_{i=1}^{n} x_i y_i$
    \end{itemize}
    
    \item \textbf{Chuẩn Euclidean:} Chuẩn Euclidean (hay chuẩn $L^2$) của một vector $\mathbf{x}$ được định nghĩa là:
    \begin{equation}
        \|\mathbf{x}\| = \sqrt{\sum_{i=1}^{n} x_i^2} = \sqrt{\langle \mathbf{x}, \mathbf{x} \rangle}
    \end{equation}
    
    \item \textbf{Tập hợp chỉ số:} Ký hiệu $[n]$ biểu thị tập hợp các số nguyên dương từ 1 đến $n$:
    \begin{equation}
        [n] = \{1, 2, \ldots, n\}
    \end{equation}
    
    \item \textbf{Bài toán hồi quy và hàm loss:} Xét một mạng nơ-ron với tham số $\theta$ thực hiện ánh xạ đầu vào $\mathbf{x}$ sang đầu ra $F(\mathbf{x}; \theta)$. Cho tập dữ liệu huấn luyện $\{(\mathbf{x}_i, \mathbf{y}_i)\}_{i=1}^{N}$ với $N$ mẫu, mục tiêu huấn luyện là tối thiểu hóa hàm loss Mean Squared Error (MSE):
    \begin{equation}
        \mathcal{L}(\theta) = \frac{1}{2} \|F(\mathbf{x}; \theta) - \mathbf{y}\|^2 = \frac{1}{2} \sum_{i=1}^{N} \|F(\mathbf{x}_i; \theta) - \mathbf{y}_i\|^2
    \end{equation}
    Trong đó $\mathbf{y}$ biểu thị vector chứa tất cả các nhãn đầu ra mục tiêu.
    
    \begin{itemize}
      \item \textbf{Tại sao chọn MSE?}
      \begin{itemize}
          \item Đạo hàm $\frac{\partial \mathcal{L}}{\partial F} = F - \mathbf{y}$ có dạng tuyến tính, cho phép gradient descent hội tụ ổn định.
          \item Trong chế độ NTK, MSE loss dẫn đến nghiệm dạng tường minh theo hồi quy hạt nhân, tạo nền tảng cho phân tích lý thuyết chính xác.
      \end{itemize}
      $\Rightarrow$ Nhờ đó, ta có thể chứng minh rằng dấu của đầu ra dự đoán luôn khớp với nhãn mục tiêu (tức là nếu $y = +1$ thì $\hat{y} > 0$, và ngược lại). Điều này đảm bảo rằng sau khi làm tròn, kết quả sẽ chính xác tuyệt đối - đây là nền tảng cho khái niệm \textit{Khả năng học chính xác} được trình bày trong Định lý 5.1 (Chương 3).
    \end{itemize}
\end{enumerate}

Lệnh \textbf{SBN} có dạng "\textit{trừ và nếu kết quả âm thì nhảy}", nên thỏa cả cập nhật trạng thái, điều kiện, và điều khiển luồng; do đó có thể đạt tính đầy đủ Turing khi được lặp lại/ghép chuỗi phù hợp \cite{minsky_1967,hopcroft_ullman_1979}.

Tác giả đưa SBN vào để nhấn mạnh rằng khung khớp mẫu và phân tích NTK không chỉ nhắm tới các phép toán số học cụ thể, mà còn đóng vai trò \textbf{nền tảng lý thuyết} hướng tới việc học và thực thi \textbf{mọi thuật toán} \cite{de_luca_2025}.

\subsection{Neural Tangent Kernel (NTK)}

NTK là một công cụ lý thuyết mạnh mẽ để phân tích động lực học huấn luyện của mạng nơ-ron \cite{jacot_2018}. Ý tưởng cốt lõi là mô tả sự thay đổi của đầu ra mạng theo thời gian huấn luyện thông qua một kernel cố định.

\subsubsection{Ma trận Jacobian}

Xét mạng nơ-ron với đầu ra $F(\mathbf{x}; \theta) \in \mathbb{R}^k$ phụ thuộc vào vector tham số $\theta \in \mathbb{R}^p$. Ma trận Jacobian của đầu ra theo tham số được định nghĩa là:
\begin{equation}
    J(\theta) = \frac{\partial F(\mathbf{x}; \theta)}{\partial \theta} \in \mathbb{R}^{Nk \times p}
\end{equation}

trong đó $N$ là số lượng mẫu dữ liệu, $k$ là chiều của đầu ra, và $p$ là số lượng tham số. Mỗi hàng của ma trận Jacobian biểu thị gradient của một thành phần đầu ra theo toàn bộ tham số.

\subsubsection{Định nghĩa NTK thực nghiệm}

Neural Tangent Kernel thực nghiệm được định nghĩa thông qua tích của ma trận Jacobian với chuyển vị của nó:
\begin{equation}
    K(\theta) = J(\theta) J(\theta)^\top \in \mathbb{R}^{Nk \times Nk}
\end{equation}

Phần tử $(i,j)$ của ma trận kernel này đo lường mức độ tương quan giữa gradient của đầu ra thứ $i$ và đầu ra thứ $j$ trong không gian tham số. Cụ thể hơn, với hai điểm dữ liệu $\mathbf{x}$ và $\mathbf{x}'$, giá trị kernel được tính như sau:
\begin{equation}
    K(\mathbf{x}, \mathbf{x}'; \theta) = \left\langle \frac{\partial F(\mathbf{x}; \theta)}{\partial \theta}, \frac{\partial F(\mathbf{x}'; \theta)}{\partial \theta} \right\rangle
\end{equation}

\subsubsection{Ý nghĩa vật lý của NTK}

NTK cho phép hiểu rõ cách mạng nơ-ron \textit{học} từ dữ liệu. Khi huấn luyện bằng gradient descent, sự thay đổi của đầu ra mạng tại một điểm $\mathbf{x}$ không chỉ phụ thuộc vào gradient tại điểm đó mà còn phụ thuộc vào gradient tại tất cả các điểm dữ liệu khác thông qua kernel. Đây chính là cơ chế \textit{học chuyển giao} nội tại trong mạng nơ-ron: thông tin học được từ một mẫu có thể ảnh hưởng đến dự đoán tại các mẫu khác.

\subsection{Giới hạn chiều rộng vô hạn và Động lực học}

Một trong những đóng góp quan trọng nhất của lý thuyết NTK là khám phá hành vi của mạng nơ-ron khi chiều rộng lớp ẩn tiến tới vô hạn. Trong giới hạn này, việc huấn luyện mạng nơ-ron trở nên đơn giản hóa đáng kể và có thể được phân tích chính xác.

\subsubsection{Chế độ huấn luyện lười}

Khi chiều rộng lớp ẩn $m \rightarrow \infty$, mạng nơ-ron hoạt động trong chế độ được gọi là \textit{huấn luyện lười}. Trong chế độ này, các tham số của mạng thay đổi cực nhỏ quanh điểm khởi tạo trong suốt quá trình huấn luyện:
\begin{equation}
    \|\theta(t) - \theta(0)\| \approx O\left(\frac{1}{\sqrt{m}}\right) \rightarrow 0 \quad \text{khi } m \rightarrow \infty
\end{equation}

Điều này có vẻ nghịch lý: làm sao mạng có thể học được nếu tham số hầu như không thay đổi? Câu trả lời nằm ở việc với số lượng tham số cực lớn, ngay cả những thay đổi nhỏ trong mỗi tham số cũng đủ để tạo ra thay đổi đáng kể trong đầu ra. Đây chính là sức mạnh của việc có nhiều chiều trong không gian tham số.

\subsubsection{Sự ổn định của Kernel (Kernel đóng băng)}

Một hệ quả quan trọng của huấn luyện lười là NTK trở nên ổn định (hay đóng băng) trong suốt quá trình huấn luyện. Cụ thể, khi $m \rightarrow \infty$:
\begin{equation}
    K_m(\theta(t)) \rightarrow \Theta(\mathbf{x}, \mathbf{x}') \quad \text{với mọi } t \geq 0
\end{equation}

trong đó $\Theta(\mathbf{x}, \mathbf{x}')$ là \textit{NTK với chiều rộng lớp ẩn vô hạn}, một kernel xác định và cố định chỉ phụ thuộc vào cặp đầu vào $(\mathbf{x}, \mathbf{x}')$ và kiến trúc mạng, không phụ thuộc vào quá trình huấn luyện. Sự ổn định này là chìa khóa cho phép phân tích chính xác động lực học huấn luyện.

\subsubsection{Phương trình vi phân cho đầu ra mạng}

Xét việc huấn luyện mạng bằng gradient flow (gradient descent với learning rate vô cùng nhỏ). Động lực học của tham số được mô tả bởi:
\begin{equation}
    \frac{d\theta}{dt} = -\nabla_\theta \mathcal{L}(\theta) = -J(\theta)^\top (F(\mathbf{x}; \theta) - \mathbf{y})
\end{equation}

Từ đây, có thể suy ra phương trình vi phân cho đầu ra mạng. Đặt $f(t) = F(\mathbf{x}; \theta(t))$ là vector đầu ra tại thời điểm $t$, ta có:
\begin{equation}
    \frac{df}{dt} = J(\theta) \frac{d\theta}{dt} = -J(\theta) J(\theta)^\top (f(t) - \mathbf{y}) = -K(\theta(t))(f(t) - \mathbf{y})
\end{equation}

Trong giới hạn chiều rộng vô hạn, với $K(\theta(t)) \approx \Theta$ cố định, phương trình trên trở thành phương trình vi phân tuyến tính:
\begin{equation}
    \frac{df}{dt} = -\Theta(f(t) - \mathbf{y})
\end{equation}

Nghiệm của phương trình này có dạng giải tích tường minh:
\begin{equation}
    f(t) = \mathbf{y} + e^{-\Theta t}(f(0) - \mathbf{y})
\end{equation}

trong đó $e^{-\Theta t}$ là hàm mũ của ma trận. Khi $t \rightarrow \infty$, nếu $\Theta$ là ma trận xác định dương, ta có $f(t) \rightarrow \mathbf{y}$, nghĩa là mạng học được chính xác nhãn đầu ra.

\subsubsection{Tương đương với Hồi quy hạt nhân}

Kết quả quan trọng nhất của phân tích trên là việc huấn luyện mạng nơ-ron trong giới hạn chiều rộng vô hạn hoàn toàn tương đương với bài toán \textit{Hồi quy hạt nhân} cổ điển. Cho một điểm test mới $\mathbf{x}^*$, dự đoán của mạng sau khi huấn luyện hội tụ là:
\begin{equation}
    f_\infty(\mathbf{x}^*) = \Theta(\mathbf{x}^*, \mathcal{X}) \Theta(\mathcal{X}, \mathcal{X})^{-1} \mathbf{y}
\end{equation}

trong đó:
\begin{quote}
\begin{itemize}
    \item $\mathcal{X} = \{\mathbf{x}_1, \ldots, \mathbf{x}_N\}$ là tập các điểm huấn luyện
    \item $\Theta(\mathbf{x}^*, \mathcal{X}) \in \mathbb{R}^{1 \times N}$ là vector kernel giữa điểm test và các điểm huấn luyện
    \item $\Theta(\mathcal{X}, \mathcal{X}) \in \mathbb{R}^{N \times N}$ là ma trận kernel (ma trận Gram) trên tập huấn luyện
\end{itemize}
\end{quote}

$\Rightarrow$ Công thức này chính là nghiệm của bài toán Hồi quy hạt nhân tiêu chuẩn. Điều này cho thấy mạng nơ-ron sâu, trong giới hạn chiều rộng vô hạn, hoạt động như một phương pháp kernel cổ điển với một kernel đặc biệt được xác định bởi kiến trúc mạng.

\subsection{Kiến trúc mạng nơ-ron hai lớp}

Nghiên cứu tập trung vào mạng nơ-ron hai lớp với kiến trúc đơn giản nhưng đủ mạnh để phân tích lý thuyết.

\subsubsection{Cấu trúc mạng}

Mạng nơ-ron hai lớp được sử dụng có cấu trúc như sau:
\begin{quote}
\begin{itemize}
    \item \textbf{Lớp đầu vào:} Nhận vector $\mathbf{x} \in \mathbb{R}^d$
    \item \textbf{Lớp ẩn:} Gồm $m$ nơ-ron với hàm kích hoạt ReLU, không sử dụng hệ số chệch
    \item \textbf{Lớp đầu ra:} Ánh xạ tuyến tính từ $m$ nơ-ron ẩn sang $k$ đầu ra
\end{itemize}
\end{quote}

Công thức toán học của mạng được viết là:
\begin{equation}
    F(\mathbf{x}) = W^2 \cdot \text{ReLU}(W^1 \mathbf{x})
\end{equation}

trong đó:
\begin{quote}
\begin{itemize}
    \item $W^1 \in \mathbb{R}^{m \times d}$ là ma trận trọng số lớp thứ nhất
    \item $W^2 \in \mathbb{R}^{k \times m}$ là ma trận trọng số lớp thứ hai
    \item $\text{ReLU}(z) = \max(0, z)$ là hàm kích hoạt ReLU được áp dụng trên từng phần tử
\end{itemize}
\end{quote}

Việc không sử dụng hệ số chệch đơn giản hóa phân tích lý thuyết mà không làm mất tính tổng quát về mặt biểu diễn, vì hệ số chệch có thể được mô phỏng bằng cách thêm một chiều hằng số vào đầu vào.

\subsubsection{Tham số hóa NTK}

Để đảm bảo tính hội tụ trong giới hạn chiều rộng vô hạn, các trọng số được khởi tạo theo chuẩn \textit{Tham số hóa NTK}. Cụ thể:
\begin{quote}
\begin{itemize}
    \item Các phần tử của $W^1$ được khởi tạo i.i.d. từ phân phối $\mathcal{N}(0, 1)$
    \item Các phần tử của $W^2$ được khởi tạo i.i.d. từ phân phối $\mathcal{N}(0, 1/m)$
\end{itemize}
\end{quote}

Hệ số $1/m$ trong phương sai của $W^2$ là then chốt. Nó đảm bảo rằng đầu ra của mạng có phương sai $O(1)$ khi $m \rightarrow \infty$:
\begin{equation}
    \text{Var}[F(\mathbf{x})] = \text{Var}\left[\sum_{j=1}^{m} W^2_{ij} \cdot \text{ReLU}(W^1_j \mathbf{x})\right] \approx m \cdot \frac{1}{m} \cdot O(1) = O(1)
\end{equation}

Nếu không có hệ số này, đầu ra sẽ có phương sai tăng theo $m$, làm mất tính ổn định khi $m \rightarrow \infty$.

